\documentclass{beamer}
\usetheme{Madrid}
\usecolortheme{default}
\usepackage[utf8]{inputenc}
\usepackage{amsmath}
\usepackage{amsfonts}
\usepackage{amssymb}
\usepackage{tikz}

\title{CRISC Exam - Hard Questions with Explanations}
\subtitle{From Open Source CRISC Exam Guide}
\author{Compiled from GTPs and AI}
\date{\today}

\begin{document}
	
	\begin{frame}
		\titlepage
	\end{frame}
	
	\begin{frame}{Outline}
		\tableofcontents
	\end{frame}
	
	% ================ SECTION 1: DOMAIN 1 - GOVERNANCE ================
	
	\section{Domain 1: Governance}
	
	\begin{frame}{Q1: Organizational Governance}
		\begin{block}{Question}
			The purpose of having a governance structure is to:
		\end{block}
		\begin{enumerate}[(A)]
			\item Manage day-to-day operations
			\item Have accountability for business alignment and operations
			\item Implement security controls
			\item Conduct risk assessments
		\end{enumerate}
	\end{frame}
	
	\begin{frame}{Q1: Explanation}
		\begin{block}{Correct Answer: B}
			The purpose of having a governance structure is to have accountability for the business alignment and day-to-day operations of the organization.
		\end{block}
		\begin{block}{Option Analysis}
			\begin{itemize}
				\item \textbf{A - Incorrect:} Management handles day-to-day operations, not governance
				\item \textbf{B - Correct:} Governance establishes accountability for business alignment
				\item \textbf{C - Incorrect:} Controls are implemented at management level
				\item \textbf{D - Incorrect:} Risk assessments are operational activities
			\end{itemize}
		\end{block}
	\end{frame}
	
	\begin{frame}{Q2: Risk Management First Step}
		\begin{block}{Question}
			Which of the following is the first step in IT risk management?
		\end{block}
		\begin{enumerate}[(A)]
			\item IT risk assessment
			\item Risk response and mitigation
			\item IT risk identification
			\item Risk reporting
		\end{enumerate}
	\end{frame}
	
	\begin{frame}{Q2: Explanation}
		\begin{block}{Correct Answer: C}
			IT risk identification is the first step of IT risk management. You cannot protect what you do not know exists.
		\end{block}
		\begin{block}{Option Analysis}
			\begin{itemize}
				\item \textbf{A - Incorrect:} Assessment comes after identification
				\item \textbf{B - Incorrect:} Response comes after assessment
				\item \textbf{C - Correct:} Identification is the first step
				\item \textbf{D - Incorrect:} Reporting is the final step
			\end{itemize}
		\end{block}
	\end{frame}
	
	\begin{frame}{Q3: Oversight Responsibility}
		\begin{block}{Question}
			Oversight and direction are the responsibility of who?
		\end{block}
		\begin{enumerate}[(A)]
			\item Management
			\item Senior staff
			\item IT risk manager
			\item Board of directors
		\end{enumerate}
	\end{frame}
	
	\begin{frame}{Q3: Explanation}
		\begin{block}{Correct Answer: D}
			Oversight and direction are the responsibility of the board of directors.
		\end{block}
		\begin{block}{Option Analysis}
			\begin{itemize}
				\item \textbf{A - Incorrect:} Management executes, doesn't provide oversight
				\item \textbf{B - Incorrect:} Senior staff don't provide governance oversight
				\item \textbf{C - Incorrect:} Risk manager executes risk activities
				\item \textbf{D - Correct:} Board provides oversight and direction
			\end{itemize}
		\end{block}
	\end{frame}
	
	\begin{frame}{Q4: Policy Documentation Hierarchy}
		\begin{block}{Question}
			The high-level intent of an organization's security practices is best established by which of the following?
		\end{block}
		\begin{enumerate}[(A)]
			\item Procedures
			\item Policies
			\item Guidelines
			\item Processes
		\end{enumerate}
	\end{frame}
	
	\begin{frame}{Q4: Explanation}
		\begin{block}{Correct Answer: B}
			Policies are the high-level intent of organizations for security practices.
		\end{block}
		\begin{block}{Option Analysis}
			\begin{itemize}
				\item \textbf{A - Incorrect:} Procedures are detailed steps
				\item \textbf{B - Correct:} Policies express management intent
				\item \textbf{C - Incorrect:} Guidelines are recommendations
				\item \textbf{D - Incorrect:} Processes are operational
			\end{itemize}
		\end{block}
	\end{frame}
	
	\begin{frame}{Q5: RACI Model}
		\begin{block}{Question}
			Which group should be reached out to for expert guidance on the working group?
		\end{block}
		\begin{enumerate}[(A)]
			\item Responsible
			\item Accountable
			\item Consulted
			\item Informed
		\end{enumerate}
	\end{frame}
	
	\begin{frame}{Q5: Explanation}
		\begin{block}{Correct Answer: C}
			Consulted working groups are the subject matter experts for guidance on initiatives in their domain of expertise.
		\end{block}
		\begin{block}{Option Analysis}
			\begin{itemize}
				\item \textbf{A - Incorrect:} Responsible does the work
				\item \textbf{B - Incorrect:} Accountable provides resources
				\item \textbf{C - Correct:} Consulted are subject matter experts
				\item \textbf{D - Incorrect:} Informed are kept updated
			\end{itemize}
		\end{block}
	\end{frame}
	
	\begin{frame}{Q6: Risk Culture Focus}
		\begin{block}{Question}
			An IT risk manager recently joined an organization with past security incidents. What should be the primary focus in the first few months?
		\end{block}
		\begin{enumerate}[(A)]
			\item Cite individuals who caused incidents
			\item Understand the organization's culture toward risk
			\item Draft strategy for IT risk management
			\item Complete all security awareness training
		\end{enumerate}
	\end{frame}
	
	\begin{frame}{Q6: Explanation}
		\begin{block}{Correct Answer: B}
			The risk manager needs to understand the culture of the organization before drafting strategy.
		\end{block}
		\begin{block}{Option Analysis}
			\begin{itemize}
				\item \textbf{A - Incorrect:} Blaming individuals is counterproductive
				\item \textbf{B - Correct:} Culture understanding is foundational
				\item \textbf{C - Incorrect:} Strategy comes after understanding culture
				\item \textbf{D - Incorrect:} Training is secondary to culture assessment
			\end{itemize}
		\end{block}
	\end{frame}
	
	\section{Domain 1: 3 Lines of Defense}
	
	\begin{frame}{Q7: 3LoD Risk Monitoring}
		\begin{block}{Question}
			In the 3LoD model, which LoD is responsible for risk monitoring and oversight?
		\end{block}
		\begin{enumerate}[(A)]
			\item First LoD
			\item Second LoD
			\item Third LoD
			\item All of the above
		\end{enumerate}
	\end{frame}
	
	\begin{frame}{Q7: Explanation}
		\begin{block}{Correct Answer: B}
			The second LoD is responsible for risk monitoring and oversight.
		\end{block}
		\begin{block}{Option Analysis}
			\begin{itemize}
				\item \textbf{A - Incorrect:} First LoD is operational management
				\item \textbf{B - Correct:} Second LoD provides oversight and monitoring
				\item \textbf{C - Incorrect:} Third LoD provides independent assurance
				\item \textbf{D - Incorrect:} Only second LoD has this responsibility
			\end{itemize}
		\end{block}
	\end{frame}
	
	\begin{frame}{Q8: Third LoD Responsibility}
		\begin{block}{Question}
			What is the primary responsibility of the third LoD?
		\end{block}
		\begin{enumerate}[(A)]
			\item Policy and procedure development
			\item Provide independent assurance of controls
			\item Perform periodic user reviews
			\item Restrict least privileged roles
		\end{enumerate}
	\end{frame}
	
	\begin{frame}{Q8: Explanation}
		\begin{block}{Correct Answer: B}
			The third LoD is primarily responsible for independent assurance of controls through internal/external audit.
		\end{block}
		\begin{block}{Option Analysis}
			\begin{itemize}
				\item \textbf{A - Incorrect:} Policy development is second LoD
				\item \textbf{B - Correct:} Independence is key for third LoD
				\item \textbf{C - Incorrect:} User reviews are operational (first LoD)
				\item \textbf{D - Incorrect:} Access control is operational
			\end{itemize}
		\end{block}
	\end{frame}
	
	\begin{frame}{Q9: Risk Appetite Definition}
		\begin{block}{Question}
			The amount of risk an organization is willing to accept is known as:
		\end{block}
		\begin{enumerate}[(A)]
			\item Risk tolerance
			\item Risk capacity
			\item Risk profile
			\item Risk appetite
		\end{enumerate}
	\end{frame}
	
	\begin{frame}{Q9: Explanation}
		\begin{block}{Correct Answer: D}
			Risk appetite is the amount of risk an organization is willing to accept to achieve its objectives.
		\end{block}
		\begin{block}{Option Analysis}
			\begin{itemize}
				\item \textbf{A - Incorrect:} Tolerance is acceptable variation
				\item \textbf{B - Incorrect:} Capacity is maximum affordable risk
				\item \textbf{C - Incorrect:} Profile is overall exposure
				\item \textbf{D - Correct:} Appetite is willingness to accept
			\end{itemize}
		\end{block}
	\end{frame}
	
	\begin{frame}{Q10: Risk Acceptance Reason}
		\begin{block}{Question}
			The BEST reason for a Business Process Owner to accept risk is:
		\end{block}
		\begin{enumerate}[(A)]
			\item Difficulty to implement controls
			\item Unavailability of resources
			\item Cost of control outweighs cost of assets
			\item Budgetary constraints
		\end{enumerate}
	\end{frame}
	
	\begin{frame}{Q10: Explanation}
		\begin{block}{Correct Answer: C}
			Cost of control implementation should not outweigh the cost of assets.
		\end{block}
		\begin{block}{Option Analysis}
			\begin{itemize}
				\item \textbf{A - Incorrect:} Difficulty is not a valid reason
				\item \textbf{B - Incorrect:} Resource availability is temporary
				\item \textbf{C - Correct:} Cost-benefit analysis is fundamental
				\item \textbf{D - Incorrect:} Budget constraints are not primary
			\end{itemize}
		\end{block}
	\end{frame}
	
	\begin{frame}{Q11: Risk Tolerance vs Capacity}
		\begin{block}{Question}
			Which statement is correct?
		\end{block}
		\begin{enumerate}[(A)]
			\item Breaching risk tolerance threatens existence
			\item Breaching risk capacity threatens existence
			\item Risk tolerance and capacity are not related
			\item Risk tolerance and capacity are the same
		\end{enumerate}
	\end{frame}
	
	\begin{frame}{Q11: Explanation}
		\begin{block}{Correct Answer: B}
			Breaching risk capacity could threaten an organization's existence.
		\end{block}
		\begin{block}{Option Analysis}
			\begin{itemize}
				\item \textbf{A - Incorrect:} Tolerance breach is manageable
				\item \textbf{B - Correct:} Capacity breach threatens existence
				\item \textbf{C - Incorrect:} They are related
				\item \textbf{D - Incorrect:} They are different concepts
			\end{itemize}
		\end{block}
	\end{frame}
	
	\section{Domain 1: Legal and Ethics}
	
	\begin{frame}{Q12: HIPAA}
		\begin{block}{Question}
			Which federal law provides guidance on protecting sensitive health information?
		\end{block}
		\begin{enumerate}[(A)]
			\item CCPA
			\item HIPAA
			\item FFIEC
			\item GDPR
		\end{enumerate}
	\end{frame}
	
	\begin{frame}{Q12: Explanation}
		\begin{block}{Correct Answer: B}
			HIPAA is the federal law protecting sensitive health information.
		\end{block}
		\begin{block}{Option Analysis}
			\begin{itemize}
				\item \textbf{A - Incorrect:} CCPA is California privacy law
				\item \textbf{B - Correct:} HIPAA is health information law
				\item \textbf{C - Incorrect:} FFIEC is banking regulation
				\item \textbf{D - Incorrect:} GDPR is EU regulation
			\end{itemize}
		\end{block}
	\end{frame}
	
	\begin{frame}{Q13: CCPA Beneficiaries}
		\begin{block}{Question}
			Which group is a beneficiary of CCPA law?
		\end{block}
		\begin{enumerate}[(A)]
			\item US residents
			\item EU residents
			\item California residents
			\item Canada residents
		\end{enumerate}
	\end{frame}
	
	\begin{frame}{Q13: Explanation}
		\begin{block}{Correct Answer: C}
			CCPA protects California residents' personal information.
		\end{block}
		\begin{block}{Option Analysis}
			\begin{itemize}
				\item \textbf{A - Incorrect:} US residents is too broad
				\item \textbf{B - Incorrect:} EU residents are covered by GDPR
				\item \textbf{C - Correct:} California residents
				\item \textbf{D - Incorrect:} Canada residents have PIPEDA
			\end{itemize}
		\end{block}
	\end{frame}
	
	\begin{frame}{Q14: HIPAA Breach Notification}
		\begin{block}{Question}
			According to HIPAA Breach Notification Rule, when must a CE report a breach to OCR?
		\end{block}
		\begin{enumerate}[(A)]
			\item Breach of more than 100 California residents
			\item IT incident in the CE
			\item IT incident in business associate
			\item Breach of health information of more than 500 individuals
		\end{enumerate}
	\end{frame}
	
	\begin{frame}{Q14: Explanation}
		\begin{block}{Correct Answer: D}
			Breach of more than 500 individuals must be reported to OCR.
		\end{block}
		\begin{block}{Option Analysis}
			\begin{itemize}
				\item \textbf{A - Incorrect:} 100 is not the threshold
				\item \textbf{B - Incorrect:} IT incident alone is not a breach
				\item \textbf{C - Incorrect:} Business associate incident may be reported
				\item \textbf{D - Correct:} 500+ individuals requires notification
			\end{itemize}
		\end{block}
	\end{frame}
	
	\begin{frame}{Q15: ISACA Code of Ethics}
		\begin{block}{Question}
			Failure to comply with ISACA Code of Ethics could lead to:
		\end{block}
		\begin{enumerate}[(A)]
			\item Mandatory training from ISACA
			\item Additional questions in exam
			\item Investigation and disciplinary measures
			\item Immediate revocation of certification
		\end{enumerate}
	\end{frame}
	
	\begin{frame}{Q15: Explanation}
		\begin{block}{Correct Answer: C}
			ISACA will conduct investigation and take disciplinary actions upon failure to comply with Code of Ethics.
		\end{block}
		\begin{block}{Option Analysis}
			\begin{itemize}
				\item \textbf{A - Incorrect:} Training is not the only consequence
				\item \textbf{B - Incorrect:} Exam questions are not affected
				\item \textbf{C - Correct:} Investigation and disciplinary measures
				\item \textbf{D - Incorrect:} Revocation may follow investigation
			\end{itemize}
		\end{block}
	\end{frame}
	
	% ================ SECTION 2: DOMAIN 2 - IT RISK ASSESSMENT ================
	
	\section{Domain 2: Risk Management Life Cycle}
	
	\begin{frame}{Q16: IT Risk vs Enterprise Risk}
		\begin{block}{Question}
			Which depicts the correct relationship between IT risk and enterprise risk?
		\end{block}
		\begin{enumerate}[(A)]
			\item Enterprise risk is part of IT risk
			\item IT risk is part of enterprise risk
			\item These types of risk are not related
			\item Enterprise risk is independent of IT risk
		\end{enumerate}
	\end{frame}
	
	\begin{frame}{Q16: Explanation}
		\begin{block}{Correct Answer: B}
			IT risk is a subset of enterprise risk.
		\end{block}
		\begin{block}{Option Analysis}
			\begin{itemize}
				\item \textbf{A - Incorrect:} Enterprise risk is broader
				\item \textbf{B - Correct:} IT risk is part of enterprise risk
				\item \textbf{C - Incorrect:} They are related
				\item \textbf{D - Incorrect:} IT risk impacts enterprise risk
			\end{itemize}
		\end{block}
	\end{frame}
	
	\begin{frame}{Q17: Risk Identification Principle}
		\begin{block}{Question}
			Which risk management life cycle step emphasizes "You can't protect what you don't know exists"?
		\end{block}
		\begin{enumerate}[(A)]
			\item Risk and control monitoring
			\item Risk assessment
			\item Risk identification
			\item Risk categorization
		\end{enumerate}
	\end{frame}
	
	\begin{frame}{Q17: Explanation}
		\begin{block}{Correct Answer: C}
			Risk identification is the first step - you must know risks exist.
		\end{block}
		\begin{block}{Option Analysis}
			\begin{itemize}
				\item \textbf{A - Incorrect:} Monitoring happens after identification
				\item \textbf{B - Incorrect:} Assessment comes after identification
				\item \textbf{C - Correct:} Identification discovers risks
				\item \textbf{D - Incorrect:} Categorization comes after identification
			\end{itemize}
		\end{block}
	\end{frame}
	
	\begin{frame}{Q18: Formal Risk Assessment Requirement}
		\begin{block}{Question}
			Which is a formal requirement by many legal and regulatory compliance frameworks?
		\end{block}
		\begin{enumerate}[(A)]
			\item Performing vulnerability assessments
			\item Maintaining asset inventory
			\item Deploying changes without testing
			\item Conducting formal risk assessment
		\end{enumerate}
	\end{frame}
	
	\begin{frame}{Q18: Explanation}
		\begin{block}{Correct Answer: D}
			Conducting formal risk assessment is a requirement of many legal and regulatory compliance frameworks.
		\end{block}
		\begin{block}{Option Analysis}
			\begin{itemize}
				\item \textbf{A - Incorrect:} Vulnerability assessments are important but not universally mandatory
				\item \textbf{B - Incorrect:} Asset inventory is good practice but not always required
				\item \textbf{C - Incorrect:} This is poor practice
				\item \textbf{D - Correct:} Formal risk assessment is widely mandated
			\end{itemize}
		\end{block}
	\end{frame}
	
	\begin{frame}{Q19: Event vs Incident vs Breach}
		\begin{block}{Question}
			A healthcare employee accidentally sent more than 5,000 patient records in response to a phishing email. This is a:
		\end{block}
		\begin{enumerate}[(A)]
			\item Issue
			\item Event
			\item Incident
			\item Breach
		\end{enumerate}
	\end{frame}
	
	\begin{frame}{Q19: Explanation}
		\begin{block}{Correct Answer: D}
			This is a breach - unauthorized disclosure of PHI affecting more than 500 individuals.
		\end{block}
		\begin{block}{Option Analysis}
			\begin{itemize}
				\item \textbf{A - Incorrect:} Issue is a condition, not an occurrence
				\item \textbf{B - Incorrect:} Event alone doesn't cause negative outcome
				\item \textbf{C - Incorrect:} Incident is one or more events with negative outcome
				\item \textbf{D - Correct:} Breach is unauthorized disclosure
			\end{itemize}
		\end{block}
	\end{frame}
	
	\section{Domain 2: Threat, Vulnerability, Risk}
	
	\begin{frame}{Q20: Vulnerability Definition}
		\begin{block}{Question}
			Which is considered a weakness that could be exploited by a malicious actor?
		\end{block}
		\begin{enumerate}[(A)]
			\item Threat
			\item Threat actor
			\item Vulnerability
			\item Risk
		\end{enumerate}
	\end{frame}
	
	\begin{frame}{Q20: Explanation}
		\begin{block}{Correct Answer: C}
			Vulnerabilities are weaknesses that can be exploited by malicious actors.
		\end{block}
		\begin{block}{Option Analysis}
			\begin{itemize}
				\item \textbf{A - Incorrect:} Threat is potential cause
				\item \textbf{B - Incorrect:} Threat actor executes the threat
				\item \textbf{C - Correct:} Vulnerability is a weakness
				\item \textbf{D - Incorrect:} Risk is the result of exploitation
			\end{itemize}
		\end{block}
	\end{frame}
	
	\begin{frame}{Q21: Threat Modeling Phase}
		\begin{block}{Question}
			"What could go wrong?" corresponds to which phase of threat modeling?
		\end{block}
		\begin{enumerate}[(A)]
			\item Model
			\item Identify
			\item Mitigate
			\item Validate
		\end{enumerate}
	\end{frame}
	
	\begin{frame}{Q21: Explanation}
		\begin{block}{Correct Answer: B}
			Identify phase requires asking "What could go wrong?"
		\end{block}
		\begin{block}{Option Analysis}
			\begin{itemize}
				\item \textbf{A - Incorrect:} Model is "What are we building?"
				\item \textbf{B - Correct:} Identify is "What could go wrong?"
				\item \textbf{C - Incorrect:} Mitigate is "What countermeasures?"
				\item \textbf{D - Incorrect:} Validate is "Have we performed all steps?"
			\end{itemize}
		\end{block}
	\end{frame}
	
	\begin{frame}{Q22: CVSS}
		\begin{block}{Question}
			Which scoring technique is used to quantify a vulnerability?
		\end{block}
		\begin{enumerate}[(A)]
			\item CSS
			\item CSV
			\item CVSS
			\item CVVS
		\end{enumerate}
	\end{frame}
	
	\begin{frame}{Q22: Explanation}
		\begin{block}{Correct Answer: C}
			Common Vulnerability Scoring System (CVSS) is used to quantify vulnerabilities.
		\end{block}
		\begin{block}{Option Analysis}
			\begin{itemize}
				\item \textbf{A - Incorrect:} CSS is not vulnerability scoring
				\item \textbf{B - Incorrect:} CSV is file format
				\item \textbf{C - Correct:} CVSS is the standard
				\item \textbf{D - Incorrect:} CVVS is not standard
			\end{itemize}
		\end{block}
	\end{frame}
	
	\section{Domain 2: Risk Assessment Frameworks}
	
	\begin{frame}{Q23: Quantitative Risk Framework}
		\begin{block}{Question}
			Which framework is primarily used for quantitative risk management?
		\end{block}
		\begin{enumerate}[(A)]
			\item NIST 800-30
			\item FAIR
			\item ISO 27001
			\item ISO 27005
		\end{enumerate}
	\end{frame}
	
	\begin{frame}{Q23: Explanation}
		\begin{block}{Correct Answer: B}
			FAIR (Factor Analysis of Information Risk) is primarily used for quantitative risk management.
		\end{block}
		\begin{block}{Option Analysis}
			\begin{itemize}
				\item \textbf{A - Incorrect:} NIST 800-30 is qualitative
				\item \textbf{B - Correct:} FAIR is quantitative
				\item \textbf{C - Incorrect:} ISO 27001 is management standard
				\item \textbf{D - Incorrect:} ISO 27005 is qualitative
			\end{itemize}
		\end{block}
	\end{frame}
	
	\begin{frame}{Q24: Risk Assessment Approaches}
		\begin{block}{Question}
			A top-down risk assessment starts from the:
		\end{block}
		\begin{enumerate}[(A)]
			\item Team
			\item Individual
			\item Organization
			\item Department
		\end{enumerate}
	\end{frame}
	
	\begin{frame}{Q24: Explanation}
		\begin{block}{Correct Answer: C}
			Top-down risk assessment starts from the organization level.
		\end{block}
		\begin{block}{Option Analysis}
			\begin{itemize}
				\item \textbf{A - Incorrect:} Team is bottom-up
				\item \textbf{B - Incorrect:} Individual is bottom-up
				\item \textbf{C - Correct:} Organization is top-down
				\item \textbf{D - Incorrect:} Department is between top and bottom
			\end{itemize}
		\end{block}
	\end{frame}
	
	\begin{frame}{Q25: Qualitative Risk True/False}
		\begin{block}{Question}
			Which is NOT true about qualitative risk management?
		\end{block}
		\begin{enumerate}[(A)]
			\item Less expensive
			\item Subjective
			\item Requires complex computation
			\item Focused on severity
		\end{enumerate}
	\end{frame}
	
	\begin{frame}{Q25: Explanation}
		\begin{block}{Correct Answer: C}
			Qualitative risk management does NOT require complex computation.
		\end{block}
		\begin{block}{Option Analysis}
			\begin{itemize}
				\item \textbf{A - True:} Qualitative is less expensive
				\item \textbf{B - True:} Qualitative is subjective
				\item \textbf{C - Not True:} Complex computation is quantitative
				\item \textbf{D - True:} Focused on severity
			\end{itemize}
		\end{block}
	\end{frame}
	
	\begin{frame}{Q26: Supply Chain Risk Framework}
		\begin{block}{Question}
			Which NIST framework provides guidance for supply chain management?
		\end{block}
		\begin{enumerate}[(A)]
			\item NIST 800-161
			\item NIST 800-30
			\item NIST 800-57
			\item ISO 27001
		\end{enumerate}
	\end{frame}
	
	\begin{frame}{Q26: Explanation}
		\begin{block}{Correct Answer: A}
			NIST SP 800-161 provides guidance for supply chain risk management.
		\end{block}
		\begin{block}{Option Analysis}
			\begin{itemize}
				\item \textbf{A - Correct:} 800-161 is supply chain
				\item \textbf{B - Incorrect:} 800-30 is risk management
				\item \textbf{C - Incorrect:} 800-57 is key management
				\item \textbf{D - Incorrect:} ISO 27001 is management standard
			\end{itemize}
		\end{block}
	\end{frame}
	
	\begin{frame}{Q27: Bow Tie Analysis}
		\begin{block}{Question}
			Which risk assessment technique provides results by displaying links between possible causes, controls, and consequences in a diagram?
		\end{block}
		\begin{enumerate}[(A)]
			\item FAIR
			\item BTA
			\item Markov analysis
			\item Monte Carlo analysis
		\end{enumerate}
	\end{frame}
	
	\begin{frame}{Q27: Explanation}
		\begin{block}{Correct Answer: B}
			Bow Tie Analysis (BTA) displays links between causes, controls, and consequences in a diagram.
		\end{block}
		\begin{block}{Option Analysis}
			\begin{itemize}
				\item \textbf{A - Incorrect:} FAIR is quantitative taxonomy
				\item \textbf{B - Correct:} BTA is diagram-based
				\item \textbf{C - Incorrect:} Markov uses state analysis
				\item \textbf{D - Incorrect:} Monte Carlo uses simulation
			\end{itemize}
		\end{block}
	\end{frame}
	
	\begin{frame}{Q28: Risk Register Output}
		\begin{block}{Question}
			The results of a risk assessment should be summarized as a(n):
		\end{block}
		\begin{enumerate}[(A)]
			\item CAP
			\item Business continuity plan
			\item Organizational chart
			\item Risk register
		\end{enumerate}
	\end{frame}
	
	\begin{frame}{Q28: Explanation}
		\begin{block}{Correct Answer: D}
			Risk register is the immediate result of a risk assessment.
		\end{block}
		\begin{block}{Option Analysis}
			\begin{itemize}
				\item \textbf{A - Incorrect:} CAP is corrective action plan
				\item \textbf{B - Incorrect:} BCP is continuity planning
				\item \textbf{C - Incorrect:} Org chart shows structure
				\item \textbf{D - Correct:} Risk register documents risks
			\end{itemize}
		\end{block}
	\end{frame}
	
	\section{Domain 2: Business Impact Analysis}
	
	\begin{frame}{Q29: BIA vs Risk Assessment}
		\begin{block}{Question}
			Which statement is false regarding risk assessment and BIA?
		\end{block}
		\begin{enumerate}[(A)]
			\item BIA identifies RPO, RTO, MTD
			\item Risk assessment identifies mitigation plans
			\item BCP needs BIA
			\item BIA and risk assessment are the same
		\end{enumerate}
	\end{frame}
	
	\begin{frame}{Q29: Explanation}
		\begin{block}{Correct Answer: D}
			BIA and risk assessment are NOT the same.
		\end{block}
		\begin{block}{Option Analysis}
			\begin{itemize}
				\item \textbf{A - True:} BIA identifies recovery objectives
				\item \textbf{B - True:} Risk assessment identifies mitigation
				\item \textbf{C - True:} BCP needs BIA as input
				\item \textbf{D - Not True:} They are different processes
			\end{itemize}
		\end{block}
	\end{frame}
	
	\begin{frame}{Q30: RPO Direction}
		\begin{block}{Question}
			Which is backward-looking in relation to BIA?
		\end{block}
		\begin{enumerate}[(A)]
			\item RPO
			\item RTO
			\item MTD
			\item SSO
		\end{enumerate}
	\end{frame}
	
	\begin{frame}{Q30: Explanation}
		\begin{block}{Correct Answer: A}
			RPO (Recovery Point Objective) is backward-looking - it measures acceptable data loss.
		\end{block}
		\begin{block}{Option Analysis}
			\begin{itemize}
				\item \textbf{A - Correct:} RPO is backward-looking
				\item \textbf{B - Incorrect:} RTO is forward-looking
				\item \textbf{C - Incorrect:} MTD is total downtime
				\item \textbf{D - Incorrect:} SSO is single sign-on
			\end{itemize}
		\end{block}
	\end{frame}
	
	\begin{frame}{Q31: Residual Risk}
		\begin{block}{Question}
			The risk after implementing controls is called:
		\end{block}
		\begin{enumerate}[(A)]
			\item Current risk
			\item Residual risk
			\item Inherent risk
			\item Total risk
		\end{enumerate}
	\end{frame}
	
	\begin{frame}{Q31: Explanation}
		\begin{block}{Correct Answer: B}
			Residual risk = Inherent Risk - Controls. It is the risk remaining after controls are implemented.
		\end{block}
		\begin{block}{Option Analysis}
			\begin{itemize}
				\item \textbf{A - Incorrect:} Current risk is point-in-time
				\item \textbf{B - Correct:} Residual risk remains after controls
				\item \textbf{C - Incorrect:} Inherent risk is before controls
				\item \textbf{D - Incorrect:} Total risk is before controls
			\end{itemize}
		\end{block}
	\end{frame}
	
	\begin{frame}{Q32: Current Risk}
		\begin{block}{Question}
			The risk after implementing controls in a real-time scenario against threats is called:
		\end{block}
		\begin{enumerate}[(A)]
			\item Residual risk
			\item Inherent risk
			\item Current risk
			\item Total risk
		\end{enumerate}
	\end{frame}
	
	\begin{frame}{Q32: Explanation}
		\begin{block}{Correct Answer: C}
			Current risk is the residual risk in a real-time threat scenario at a specific point in time.
		\end{block}
		\begin{block}{Option Analysis}
			\begin{itemize}
				\item \textbf{A - Incorrect:} Residual is not real-time
				\item \textbf{B - Incorrect:} Inherent is before controls
				\item \textbf{C - Correct:} Current is real-time residual risk
				\item \textbf{D - Incorrect:} Total risk is before controls
			\end{itemize}
		\end{block}
	\end{frame}
	
	% ================ SECTION 3: DOMAIN 3 - RISK RESPONSE ================
	
	\section{Domain 3: Risk Response}
	
	\begin{frame}{Q33: Risk Owner Purpose}
		\begin{block}{Question}
			The primary reason for having a risk owner is:
		\end{block}
		\begin{enumerate}[(A)]
			\item Leverage their resources
			\item Ensure accountability
			\item Provide flexibility
			\item Seek assistance
		\end{enumerate}
	\end{frame}
	
	\begin{frame}{Q33: Explanation}
		\begin{block}{Correct Answer: B}
			The risk owner drives accountability in risk remediation.
		\end{block}
		\begin{block}{Option Analysis}
			\begin{itemize}
				\item \textbf{A - Incorrect:} Resources are secondary
				\item \textbf{B - Correct:} Accountability is primary
				\item \textbf{C - Incorrect:} Flexibility is not primary
				\item \textbf{D - Incorrect:} Assistance is secondary
			\end{itemize}
		\end{block}
	\end{frame}
	
	\begin{frame}{Q34: NOT a Risk Response Strategy}
		\begin{block}{Question}
			Which is NOT a risk response strategy?
		\end{block}
		\begin{enumerate}[(A)]
			\item Accept
			\item Transfer
			\item Treat
			\item Mitigate
		\end{enumerate}
	\end{frame}
	
	\begin{frame}{Q34: Explanation}
		\begin{block}{Correct Answer: C}
			"Treat" is sometimes used interchangeably but is not a formal risk response strategy. The four strategies are Accept, Transfer, Mitigate, and Avoid.
		\end{block}
		\begin{block}{Option Analysis}
			\begin{itemize}
				\item \textbf{A - Incorrect:} Accept is a strategy
				\item \textbf{B - Incorrect:} Transfer is a strategy
				\item \textbf{C - Correct:} Treat is not a formal strategy
				\item \textbf{D - Incorrect:} Mitigate is a strategy
			\end{itemize}
		\end{block}
	\end{frame}
	
	\begin{frame}{Q35: Compensating Controls}
		\begin{block}{Question}
			A risk manager determines primary controls are insufficient and implements compensating controls. This is:
		\end{block}
		\begin{enumerate}[(A)]
			\item Risk mitigation
			\item Risk acceptance
			\item Risk sharing
			\item Risk avoidance
		\end{enumerate}
	\end{frame}
	
	\begin{frame}{Q35: Explanation}
		\begin{block}{Correct Answer: A}
			Implementing additional controls is an example of risk mitigation.
		\end{block}
		\begin{block}{Option Analysis}
			\begin{itemize}
				\item \textbf{A - Correct:} Adding controls is mitigation
				\item \textbf{B - Incorrect:} Acceptance would do nothing
				\item \textbf{C - Incorrect:} Sharing would involve third party
				\item \textbf{D - Incorrect:} Avoidance would terminate activity
			\end{itemize}
		\end{block}
	\end{frame}
	
	\begin{frame}{Q36: Insurance as Risk Response}
		\begin{block}{Question}
			Senior management buys insurance for an earthquake-prone site. This is:
		\end{block}
		\begin{enumerate}[(A)]
			\item Risk mitigation
			\item Risk acceptance
			\item Risk sharing
			\item Risk avoidance
		\end{enumerate}
	\end{frame}
	
	\begin{frame}{Q36: Explanation}
		\begin{block}{Correct Answer: C}
			Insurance transfers/sharing the risk with the insurance company.
		\end{block}
		\begin{block}{Option Analysis}
			\begin{itemize}
				\item \textbf{A - Incorrect:} Mitigation would reduce risk
				\item \textbf{B - Incorrect:} Acceptance would retain risk
				\item \textbf{C - Correct:} Sharing through insurance
				\item \textbf{D - Incorrect:} Avoidance would close site
			\end{itemize}
		\end{block}
	\end{frame}
	
	\begin{frame}{Q37: Risk Avoidance}
		\begin{block}{Question}
			Management terminates onboarding a new tool as it doesn't adapt to changing requirements. This is:
		\end{block}
		\begin{enumerate}[(A)]
			\item Risk mitigation
			\item Risk acceptance
			\item Risk sharing
			\item Risk avoidance
		\end{enumerate}
	\end{frame}
	
	\begin{frame}{Q37: Explanation}
		\begin{block}{Correct Answer: D}
			Terminating a project to eliminate risk is risk avoidance.
		\end{block}
		\begin{block}{Option Analysis}
			\begin{itemize}
				\item \textbf{A - Incorrect:} Mitigation would modify the tool
				\item \textbf{B - Incorrect:} Acceptance would proceed with risk
				\item \textbf{C - Incorrect:} Sharing would involve third party
				\item \textbf{D - Correct:} Avoidance eliminates the risk
			\end{itemize}
		\end{block}
	\end{frame}
	
	\begin{frame}{Q38: Remote Work as Risk Response}
		\begin{block}{Question}
			An organization allows employees to work remotely due to geographical risks. This is:
		\end{block}
		\begin{enumerate}[(A)]
			\item Risk mitigation
			\item Risk acceptance
			\item Risk sharing
			\item Risk avoidance
		\end{enumerate}
	\end{frame}
	
	\begin{frame}{Q38: Explanation}
		\begin{block}{Correct Answer: A}
			Allowing remote work reduces the impact of geographical risks - this is mitigation.
		\end{block}
		\begin{block}{Option Analysis}
			\begin{itemize}
				\item \textbf{A - Correct:} Reducing impact through remote work
				\item \textbf{B - Incorrect:} Acceptance would do nothing
				\item \textbf{C - Incorrect:} Sharing would transfer to another
				\item \textbf{D - Incorrect:} Avoidance would close location
			\end{itemize}
		\end{block}
	\end{frame}
	
	\section{Domain 3: Third-Party Risk Management}
	
	\begin{frame}{Q39: SLA Breach}
		\begin{block}{Question}
			Which would bind a third party to provide monetary credits in case of service failure?
		\end{block}
		\begin{enumerate}[(A)]
			\item Master service agreement
			\item Service-level agreement
			\item Non-disclosure agreement
			\item External audit
		\end{enumerate}
	\end{frame}
	
	\begin{frame}{Q39: Explanation}
		\begin{block}{Correct Answer: B}
			A breach in the SLA allows the organization to demand monetary credit.
		\end{block}
		\begin{block}{Option Analysis}
			\begin{itemize}
				\item \textbf{A - Incorrect:} MSA is general framework
				\item \textbf{B - Correct:} SLA defines service levels and penalties
				\item \textbf{C - Incorrect:} NDA protects confidentiality
				\item \textbf{D - Incorrect:} Audit is verification
			\end{itemize}
		\end{block}
	\end{frame}
	
	\begin{frame}{Q40: NDA Purpose}
		\begin{block}{Question}
			Which should be signed with third party to protect intellectual property?
		\end{block}
		\begin{enumerate}[(A)]
			\item Master service agreement
			\item Service-level agreement
			\item Non-disclosure agreement
			\item External audit
		\end{enumerate}
	\end{frame}
	
	\begin{frame}{Q40: Explanation}
		\begin{block}{Correct Answer: C}
			NDAs protect intellectual property and interests of the organization.
		\end{block}
		\begin{block}{Option Analysis}
			\begin{itemize}
				\item \textbf{A - Incorrect:} MSA is general framework
				\item \textbf{B - Incorrect:} SLA defines service levels
				\item \textbf{C - Correct:} NDA protects IP and confidentiality
				\item \textbf{D - Incorrect:} Audit is verification
			\end{itemize}
		\end{block}
	\end{frame}
	
	\begin{frame}{Q41: External Audit Output}
		\begin{block}{Question}
			Which is NOT the final output of an external audit?
		\end{block}
		\begin{enumerate}[(A)]
			\item SOC 2 report
			\item ISO 27001 certification
			\item HITRUST certification
			\item Non-disclosure agreement
		\end{enumerate}
	\end{frame}
	
	\begin{frame}{Q41: Explanation}
		\begin{block}{Correct Answer: D}
			Non-disclosure agreements are signed internally, not the output of an external audit.
		\end{block}
		\begin{block}{Option Analysis}
			\begin{itemize}
				\item \textbf{A - Incorrect:} SOC 2 is audit output
				\item \textbf{B - Incorrect:} ISO 27001 is audit output
				\item \textbf{C - Incorrect:} HITRUST is audit output
				\item \textbf{D - Correct:} NDA is not an audit output
			\end{itemize}
		\end{block}
	\end{frame}
	
	\begin{frame}{Q42: Exception Review Frequency}
		\begin{block}{Question}
			The risk practitioner should review and verify granted exceptions are still required at least:
		\end{block}
		\begin{enumerate}[(A)]
			\item Weekly
			\item Monthly
			\item Quarterly
			\item Annually
		\end{enumerate}
	\end{frame}
	
	\begin{frame}{Q42: Explanation}
		\begin{block}{Correct Answer: D}
			Granted exceptions should be verified at least annually.
		\end{block}
		\begin{block}{Option Analysis}
			\begin{itemize}
				\item \textbf{A - Incorrect:} Weekly is too frequent
				\item \textbf{B - Incorrect:} Monthly is too frequent
				\item \textbf{C - Incorrect:} Quarterly is too frequent
				\item \textbf{D - Correct:} Annual review is standard
			\end{itemize}
		\end{block}
	\end{frame}
	
	\begin{frame}{Q43: CAB Composition}
		\begin{block}{Question}
			The CAB should consist of which of the following?
		\end{block}
		\begin{enumerate}[(A)]
			\item Only IT
			\item IT and engineering
			\item IT, engineering, and security
			\item All relevant stakeholders
		\end{enumerate}
	\end{frame}
	
	\begin{frame}{Q43: Explanation}
		\begin{block}{Correct Answer: D}
			CAB should consist of all relevant stakeholders.
		\end{block}
		\begin{block}{Option Analysis}
			\begin{itemize}
				\item \textbf{A - Incorrect:} IT alone is insufficient
				\item \textbf{B - Incorrect:} Engineering alone is insufficient
				\item \textbf{C - Incorrect:} Security is needed but not sufficient
				\item \textbf{D - Correct:} All relevant stakeholders needed
			\end{itemize}
		\end{block}
	\end{frame}
	
	\section{Domain 3: Control Design}
	
	\begin{frame}{Q44: Control Categories}
		\begin{block}{Question}
			Which is NOT a control category?
		\end{block}
		\begin{enumerate}[(A)]
			\item Additive
			\item Preventive
			\item Deterrent
			\item Detective
		\end{enumerate}
	\end{frame}
	
	\begin{frame}{Q44: Explanation}
		\begin{block}{Correct Answer: A}
			Additive is not a control category. The categories are Preventive, Detective, Corrective, Deterrent, Compensating.
		\end{block}
		\begin{block}{Option Analysis}
			\begin{itemize}
				\item \textbf{A - Correct:} Additive is not a control category
				\item \textbf{B - Incorrect:} Preventive is a category
				\item \textbf{C - Incorrect:} Deterrent is a category
				\item \textbf{D - Incorrect:} Detective is a category
			\end{itemize}
		\end{block}
	\end{frame}
	
	\begin{frame}{Q45: Proactive Control}
		\begin{block}{Question}
			Implementing controls proactively based on root cause analysis is an example of:
		\end{block}
		\begin{enumerate}[(A)]
			\item Preventive
			\item Detective
			\item Corrective
			\item Compensating
		\end{enumerate}
	\end{frame}
	
	\begin{frame}{Q45: Explanation}
		\begin{block}{Correct Answer: A}
			Implementing controls proactively is an example of preventive control.
		\end{block}
		\begin{block}{Option Analysis}
			\begin{itemize}
				\item \textbf{A - Correct:} Preventive stops issues before they occur
				\item \textbf{B - Incorrect:} Detective finds issues after
				\item \textbf{C - Incorrect:} Corrective fixes issues after
				\item \textbf{D - Incorrect:} Compensating compensates for weaknesses
			\end{itemize}
		\end{block}
	\end{frame}
	
	\begin{frame}{Q46: Detective Control}
		\begin{block}{Question}
			Reviewing audit logs in a SIEM tool is an example of:
		\end{block}
		\begin{enumerate}[(A)]
			\item Corrective control
			\item Detective control
			\item Compensating control
			\item Deterrent control
		\end{enumerate}
	\end{frame}
	
	\begin{frame}{Q46: Explanation}
		\begin{block}{Correct Answer: B}
			Reviewing audit logs is detective control - it identifies issues after they occur.
		\end{block}
		\begin{block}{Option Analysis}
			\begin{itemize}
				\item \textbf{A - Incorrect:} Corrective fixes issues
				\item \textbf{B - Correct:} Detective identifies issues
				\item \textbf{C - Incorrect:} Compensating compensates
				\item \textbf{D - Incorrect:} Deterrent discourages actions
			\end{itemize}
		\end{block}
	\end{frame}
	
	\begin{frame}{Q47: Deterrent Control}
		\begin{block}{Question}
			Installing security cameras to secure a data center is:
		\end{block}
		\begin{enumerate}[(A)]
			\item Corrective control
			\item Detective control
			\item Compensating control
			\item Deterrent control
		\end{enumerate}
	\end{frame}
	
	\begin{frame}{Q47: Explanation}
		\begin{block}{Correct Answer: D}
			Security cameras are a deterrent control - they discourage malicious activity.
		\end{block}
		\begin{block}{Option Analysis}
			\begin{itemize}
				\item \textbf{A - Incorrect:} Corrective would fix issues
				\item \textbf{B - Incorrect:} Detective would identify issues
				\item \textbf{C - Incorrect:} Compensating would compensate
				\item \textbf{D - Correct:} Deterrent discourages activity
			\end{itemize}
		\end{block}
	\end{frame}
	
	\begin{frame}{Q48: Compensating Control}
		\begin{block}{Question}
			A risk manager introduced additional oversight of an accounting system where logical controls are not implemented. This is:
		\end{block}
		\begin{enumerate}[(A)]
			\item Corrective control
			\item Detective control
			\item Compensating control
			\item Deterrent control
		\end{enumerate}
	\end{frame}
	
	\begin{frame}{Q48: Explanation}
		\begin{block}{Correct Answer: C}
			Additional oversight compensating for missing logical controls is a compensating control.
		\end{block}
		\begin{block}{Option Analysis}
			\begin{itemize}
				\item \textbf{A - Incorrect:} Corrective would fix the issue
				\item \textbf{B - Incorrect:} Detective would identify
				\item \textbf{C - Correct:} Compensating compensates for weakness
				\item \textbf{D - Incorrect:} Deterrent would discourage
			\end{itemize}
		\end{block}
	\end{frame}
	
	\begin{frame}{Q49: Corrective Control}
		\begin{block}{Question}
			An IT system failed during changes and lost data. Risk manager advises restoring from backups. This is:
		\end{block}
		\begin{enumerate}[(A)]
			\item Corrective control
			\item Detective control
			\item Compensating control
			\item Deterrent control
		\end{enumerate}
	\end{frame}
	
	\begin{frame}{Q49: Explanation}
		\begin{block}{Correct Answer: A}
			Restoring from backups to fix an issue is a corrective control.
		\end{block}
		\begin{block}{Option Analysis}
			\begin{itemize}
				\item \textbf{A - Correct:} Corrective fixes the issue
				\item \textbf{B - Incorrect:} Detective would identify
				\item \textbf{C - Incorrect:} Compensating would compensate
				\item \textbf{D - Incorrect:} Deterrent would discourage
			\end{itemize}
		\end{block}
	\end{frame}
	
	\begin{frame}{Q50: Post-Incident Controls}
		\begin{block}{Question}
			Which control categories are used after the incident has happened?
		\end{block}
		\begin{enumerate}[(A)]
			\item Corrective, compensating
			\item Corrective, deterrent
			\item Corrective, detective
			\item Deterrent, preventive
		\end{enumerate}
	\end{frame}
	
	\begin{frame}{Q50: Explanation}
		\begin{block}{Correct Answer: C}
			Detective and corrective controls are used after the incident has happened.
		\end{block}
		\begin{block}{Option Analysis}
			\begin{itemize}
				\item \textbf{A - Incorrect:} Compensating is ongoing
				\item \textbf{B - Incorrect:} Deterrent is before
				\item \textbf{C - Correct:} Detective and corrective are after
				\item \textbf{D - Incorrect:} Deterrent/preventive are before
			\end{itemize}
		\end{block}
	\end{frame}
	
	\begin{frame}{Q51: Parallel Changeover}
		\begin{block}{Question}
			Risk manager overseeing critical system upgrade with no downtime. Recommendation should be:
		\end{block}
		\begin{enumerate}[(A)]
			\item Parallel
			\item Phased
			\item Abrupt
			\item Any of the above
		\end{enumerate}
	\end{frame}
	
	\begin{frame}{Q51: Explanation}
		\begin{block}{Correct Answer: A}
			Parallel changeover allows both systems to run simultaneously with no downtime.
		\end{block}
		\begin{block}{Option Analysis}
			\begin{itemize}
				\item \textbf{A - Correct:} Parallel ensures no downtime
				\item \textbf{B - Incorrect:} Phased may have downtime
				\item \textbf{C - Incorrect:} Abrupt has high risk
				\item \textbf{D - Incorrect:} Parallel is best for no downtime
			\end{itemize}
		\end{block}
	\end{frame}
	
	\begin{frame}{Q52: Phased Changeover}
		\begin{block}{Question}
			System upgrading per module with no dependencies. Most efficient method?
		\end{block}
		\begin{enumerate}[(A)]
			\item Parallel
			\item Phased
			\item Abrupt
			\item Any of the above
		\end{enumerate}
	\end{frame}
	
	\begin{frame}{Q52: Explanation}
		\begin{block}{Correct Answer: B}
			Phased changeover is most efficient when modules are independent.
		\end{block}
		\begin{block}{Option Analysis}
			\begin{itemize}
				\item \textbf{A - Incorrect:} Parallel is expensive
				\item \textbf{B - Correct:} Phased works well with independent modules
				\item \textbf{C - Incorrect:} Abrupt is too risky
				\item \textbf{D - Incorrect:} Phased is most efficient
			\end{itemize}
		\end{block}
	\end{frame}
	
	\begin{frame}{Q53: Post-Implementation Audience}
		\begin{block}{Question}
			Who should be the target audience for a post-implementation review?
		\end{block}
		\begin{enumerate}[(A)]
			\item Developers
			\item Risk manager
			\item Risk owner
			\item All major stakeholders
		\end{enumerate}
	\end{frame}
	
	\begin{frame}{Q53: Explanation}
		\begin{block}{Correct Answer: D}
			All major stakeholders should participate in post-implementation review.
		\end{block}
		\begin{block}{Option Analysis}
			\begin{itemize}
				\item \textbf{A - Incorrect:} Developers alone insufficient
				\item \textbf{B - Incorrect:} Risk manager alone insufficient
				\item \textbf{C - Incorrect:} Risk owner alone insufficient
				\item \textbf{D - Correct:} All stakeholders needed
			\end{itemize}
		\end{block}
	\end{frame}
	
	\section{Domain 3: Monitoring and Reporting}
	
	\begin{frame}{Q54: Log Aggregation Purpose}
		\begin{block}{Question}
			What is the primary purpose of log aggregation in an organization?
		\end{block}
		\begin{enumerate}[(A)]
			\item Identify and analyze important events
			\item Synchronize clocks
			\item Eliminate need for logging
			\item Provide centralized storage only
		\end{enumerate}
	\end{frame}
	
	\begin{frame}{Q54: Explanation}
		\begin{block}{Correct Answer: A}
			Log aggregation helps identify and analyze important events and potential incidents.
		\end{block}
		\begin{block}{Option Analysis}
			\begin{itemize}
				\item \textbf{A - Correct:} Identification and analysis of events
				\item \textbf{B - Incorrect:} Clock sync is separate
				\item \textbf{C - Incorrect:} Doesn't eliminate logging
				\item \textbf{D - Incorrect:} Storage is not primary purpose
			\end{itemize}
		\end{block}
	\end{frame}
	
	\begin{frame}{Q55: SIEM Tool}
		\begin{block}{Question}
			Which is a commonly used tool for log aggregation and analysis?
		\end{block}
		\begin{enumerate}[(A)]
			\item SIEM system
			\item Vulnerability assessment tool
			\item Penetration testing tool
			\item Heat map generator
		\end{enumerate}
	\end{frame}
	
	\begin{frame}{Q55: Explanation}
		\begin{block}{Correct Answer: A}
			SIEM systems are commonly used for log aggregation and analysis.
		\end{block}
		\begin{block}{Option Analysis}
			\begin{itemize}
				\item \textbf{A - Correct:} SIEM handles log aggregation
				\item \textbf{B - Incorrect:} VA tools find vulnerabilities
				\item \textbf{C - Incorrect:} Pen testing finds exploits
				\item \textbf{D - Incorrect:} Heat maps are visualization
			\end{itemize}
		\end{block}
	\end{frame}
	
	\begin{frame}{Q56: KRI Role}
		\begin{block}{Question}
			What is the role of KRIs in risk monitoring?
		\end{block}
		\begin{enumerate}[(A)]
			\item Predict risks that may breach thresholds
			\item Measure effectiveness of controls
			\item Indicate weaknesses in controls
			\item Assess control performance
		\end{enumerate}
	\end{frame}
	
	\begin{frame}{Q56: Explanation}
		\begin{block}{Correct Answer: A}
			KRIs predict risks that may breach predefined thresholds.
		\end{block}
		\begin{block}{Option Analysis}
			\begin{itemize}
				\item \textbf{A - Correct:} KRIs predict risk breaches
				\item \textbf{B - Incorrect:} KPIs measure effectiveness
				\item \textbf{C - Incorrect:} KCIs indicate weaknesses
				\item \textbf{D - Incorrect:} KPIs assess performance
			\end{itemize}
		\end{block}
	\end{frame}
	
	\begin{frame}{Q57: Dashboard Characteristics}
		\begin{block}{Question}
			Which format of risk reporting is known for flexibility in combining qualitative and quantitative metrics?
		\end{block}
		\begin{enumerate}[(A)]
			\item Executive summary
			\item Heat map
			\item Scorecard
			\item Dashboard
		\end{enumerate}
	\end{frame}
	
	\begin{frame}{Q57: Explanation}
		\begin{block}{Correct Answer: D}
			Dashboards combine qualitative and quantitative metrics with flexibility.
		\end{block}
		\begin{block}{Option Analysis}
			\begin{itemize}
				\item \textbf{A - Incorrect:} Summary is text-based
				\item \textbf{B - Incorrect:} Heat map is qualitative only
				\item \textbf{C - Incorrect:} Scorecard is qualitative
				\item \textbf{D - Correct:} Dashboard is flexible
			\end{itemize}
		\end{block}
	\end{frame}
	
	\begin{frame}{Q58: SMART Metrics}
		\begin{block}{Question}
			What are the key characteristics of a SMART metric?
		\end{block}
		\begin{enumerate}[(A)]
			\item Specific, Measurable, Attainable, Relevant, Timely
			\item Standardized, Meaningful, Achievable, Realistic, Timeless
			\item Secure, Measurable, Aligned, Reliable, Trustworthy
			\item Significant, Measurable, Accurate, Relevant, Timed
		\end{enumerate}
	\end{frame}
	
	\begin{frame}{Q58: Explanation}
		\begin{block}{Correct Answer: A}
			SMART stands for Specific, Measurable, Attainable, Relevant, Timely.
		\end{block}
		\begin{block}{Option Analysis}
			\begin{itemize}
				\item \textbf{A - Correct:} Standard SMART definition
				\item \textbf{B - Incorrect:} Not standard SMART
				\item \textbf{C - Incorrect:} Not standard SMART
				\item \textbf{D - Incorrect:} Not standard SMART
			\end{itemize}
		\end{block}
	\end{frame}
	
	\begin{frame}{Q59: Penetration Testing}
		\begin{block}{Question}
			Which technique involves simulating a real attack to identify undiscovered vulnerabilities?
		\end{block}
		\begin{enumerate}[(A)]
			\item Self-assessment
			\item Internal audit
			\item Vulnerability assessment
			\item Penetration testing
		\end{enumerate}
	\end{frame}
	
	\begin{frame}{Q59: Explanation}
		\begin{block}{Correct Answer: D}
			Penetration testing simulates real attacks to identify vulnerabilities.
		\end{block}
		\begin{block}{Option Analysis}
			\begin{itemize}
				\item \textbf{A - Incorrect:} Self-assessment is internal review
				\item \textbf{B - Incorrect:} Audit is compliance review
				\item \textbf{C - Incorrect:} VA finds known vulnerabilities
				\item \textbf{D - Correct:} Penetration test simulates attacks
			\end{itemize}
		\end{block}
	\end{frame}
	
	% ================ SECTION 4: DOMAIN 4 - IT AND SECURITY ================
	
	\section{Domain 4: Enterprise Architecture}
	
	\begin{frame}{Q60: Technology Architecture}
		\begin{block}{Question}
			Networking devices, storage, and software are components of:
		\end{block}
		\begin{enumerate}[(A)]
			\item Business architecture
			\item Technology architecture
			\item Data architecture
			\item Application architecture
		\end{enumerate}
	\end{frame}
	
	\begin{frame}{Q60: Explanation}
		\begin{block}{Correct Answer: B}
			Networking devices, storage, and software are components of technology architecture.
		\end{block}
		\begin{block}{Option Analysis}
			\begin{itemize}
				\item \textbf{A - Incorrect:} Business architecture defines processes
				\item \textbf{B - Correct:} Technology architecture covers infrastructure
				\item \textbf{C - Incorrect:} Data architecture covers data management
				\item \textbf{D - Incorrect:} Application architecture covers software solutions
			\end{itemize}
		\end{block}
	\end{frame}
	
	\begin{frame}{Q61: CMM Defined Process}
		\begin{block}{Question}
			Which is NOT a property of the Defined process in CMM?
		\end{block}
		\begin{enumerate}[(A)]
			\item Well characterized and understood
			\item Defined at organizational level
			\item Proactive
			\item Unpredictable
		\end{enumerate}
	\end{frame}
	
	\begin{frame}{Q61: Explanation}
		\begin{block}{Correct Answer: D}
			Unpredictable is not a property of Defined process in CMM.
		\end{block}
		\begin{block}{Option Analysis}
			\begin{itemize}
				\item \textbf{A - True:} Well characterized
				\item \textbf{B - True:} Defined at organizational level
				\item \textbf{C - True:} Proactive
				\item \textbf{D - Not True:} Defined processes are predictable
			\end{itemize}
		\end{block}
	\end{frame}
	
	\begin{frame}{Q62: TCP/IP Model Layers}
		\begin{block}{Question}
			Which is Layer 3 of the TCP/IP model?
		\end{block}
		\begin{enumerate}[(A)]
			\item Physical
			\item Data link
			\item Network
			\item Application
		\end{enumerate}
	\end{frame}
	
	\begin{frame}{Q62: Explanation}
		\begin{block}{Correct Answer: C}
			Network layer is Layer 3 of the TCP/IP model.
		\end{block}
		\begin{block}{Option Analysis}
			\begin{itemize}
				\item \textbf{A - Incorrect:} Physical is Layer 1
				\item \textbf{B - Incorrect:} Data link is Layer 2
				\item \textbf{C - Correct:} Network is Layer 3
				\item \textbf{D - Incorrect:} Application is Layer 5
			\end{itemize}
		\end{block}
	\end{frame}
	
	\begin{frame}{Q63: OSI Model Session Layer}
		\begin{block}{Question}
			Which is Layer 5 of the OSI model?
		\end{block}
		\begin{enumerate}[(A)]
			\item Physical
			\item Application
			\item Session
			\item Transport
		\end{enumerate}
	\end{frame}
	
	\begin{frame}{Q63: Explanation}
		\begin{block}{Correct Answer: C}
			Session layer is Layer 5 of the OSI model.
		\end{block}
		\begin{block}{Option Analysis}
			\begin{itemize}
				\item \textbf{A - Incorrect:} Physical is Layer 1
				\item \textbf{B - Incorrect:} Application is Layer 7
				\item \textbf{C - Correct:} Session is Layer 5
				\item \textbf{D - Incorrect:} Transport is Layer 4
			\end{itemize}
		\end{block}
	\end{frame}
	
	\begin{frame}{Q64: VPN Protocol}
		\begin{block}{Question}
			Which protocol is used to implement a VPN?
		\end{block}
		\begin{enumerate}[(A)]
			\item DNS
			\item IPSec
			\item SSL
			\item SSO
		\end{enumerate}
	\end{frame}
	
	\begin{frame}{Q64: Explanation}
		\begin{block}{Correct Answer: B}
			IPSec protocol is commonly used to implement VPNs.
		\end{block}
		\begin{block}{Option Analysis}
			\begin{itemize}
				\item \textbf{A - Incorrect:} DNS is domain name service
				\item \textbf{B - Correct:} IPSec enables VPN
				\item \textbf{C - Incorrect:} SSL is for web encryption
				\item \textbf{D - Incorrect:} SSO is single sign-on
			\end{itemize}
		\end{block}
	\end{frame}
	
	\begin{frame}{Q65: IaaS}
		\begin{block}{Question}
			Which is an example of on-demand compute, network, and storage services?
		\end{block}
		\begin{enumerate}[(A)]
			\item PaaS
			\item SaaS
			\item IaaS
			\item NaaS
		\end{enumerate}
	\end{frame}
	
	\begin{frame}{Q65: Explanation}
		\begin{block}{Correct Answer: C}
			IaaS provides on-demand compute, network, and storage services.
		\end{block}
		\begin{block}{Option Analysis}
			\begin{itemize}
				\item \textbf{A - Incorrect:} PaaS provides platform
				\item \textbf{B - Incorrect:} SaaS provides software
				\item \textbf{C - Correct:} IaaS provides infrastructure
				\item \textbf{D - Incorrect:} NaaS provides network services
			\end{itemize}
		\end{block}
	\end{frame}
	
	\begin{frame}{Q66: SaaS Example}
		\begin{block}{Question}
			An online platform for video games is an example of:
		\end{block}
		\begin{enumerate}[(A)]
			\item PaaS
			\item SaaS
			\item IaaS
			\item NaaS
		\end{enumerate}
	\end{frame}
	
	\begin{frame}{Q66: Explanation}
		\begin{block}{Correct Answer: B}
			Online video game platform that lets users play without installing software is SaaS.
		\end{block}
		\begin{block}{Option Analysis}
			\begin{itemize}
				\item \textbf{A - Incorrect:} PaaS provides development platform
				\item \textbf{B - Correct:} SaaS provides ready-to-use software
				\item \textbf{C - Incorrect:} IaaS provides infrastructure
				\item \textbf{D - Incorrect:} NaaS provides network
			\end{itemize}
		\end{block}
	\end{frame}
	
	\begin{frame}{Q67: Private Cloud}
		\begin{block}{Question}
			In which deployment model are underlying cloud services NOT shared with other customers?
		\end{block}
		\begin{enumerate}[(A)]
			\item Public
			\item Private
			\item Community
			\item Hybrid
		\end{enumerate}
	\end{frame}
	
	\begin{frame}{Q67: Explanation}
		\begin{block}{Correct Answer: B}
			Private cloud resources are dedicated to a single organization.
		\end{block}
		\begin{block}{Option Analysis}
			\begin{itemize}
				\item \textbf{A - Incorrect:} Public is shared
				\item \textbf{B - Correct:} Private is dedicated
				\item \textbf{C - Incorrect:} Community is shared with similar orgs
				\item \textbf{D - Incorrect:} Hybrid combines models
			\end{itemize}
		\end{block}
	\end{frame}
	
	\begin{frame}{Q68: Community Cloud}
		\begin{block}{Question}
			In which deployment model are underlying cloud services shared with similar customers?
		\end{block}
		\begin{enumerate}[(A)]
			\item Public
			\item Private
			\item Community
			\item Hybrid
		\end{enumerate}
	\end{frame}
	
	\begin{frame}{Q68: Explanation}
		\begin{block}{Correct Answer: C}
			Community cloud lets similar organizations share resources.
		\end{block}
		\begin{block}{Option Analysis}
			\begin{itemize}
				\item \textbf{A - Incorrect:} Public is open to all
				\item \textbf{B - Incorrect:} Private is dedicated
				\item \textbf{C - Correct:} Community is for similar orgs
				\item \textbf{D - Incorrect:} Hybrid combines models
			\end{itemize}
		\end{block}
	\end{frame}
	
	\section{Domain 4: Data Life Cycle}
	
	\begin{frame}{Q69: BIA Input for BCP}
		\begin{block}{Question}
			Which acts as an input for the BCP?
		\end{block}
		\begin{enumerate}[(A)]
			\item Recovery objectives
			\item DR test
			\item BIA
			\item BC test
		\end{enumerate}
	\end{frame}
	
	\begin{frame}{Q69: Explanation}
		\begin{block}{Correct Answer: C}
			BIA is the input to creating a BCP.
		\end{block}
		\begin{block}{Option Analysis}
			\begin{itemize}
				\item \textbf{A - Incorrect:} Recovery objectives are output of BIA
				\item \textbf{B - Incorrect:} DR test validates recovery
				\item \textbf{C - Correct:} BIA inputs BCP
				\item \textbf{D - Incorrect:} BC test validates BCP
			\end{itemize}
		\end{block}
	\end{frame}
	
	\begin{frame}{Q70: RPO}
		\begin{block}{Question}
			Risk manager confirms application can't lose more than four hours of data. Which objective is this?
		\end{block}
		\begin{enumerate}[(A)]
			\item MTD
			\item RPO
			\item RTO
			\item BIA
		\end{enumerate}
	\end{frame}
	
	\begin{frame}{Q70: Explanation}
		\begin{block}{Correct Answer: B}
			RPO is the maximum acceptable data loss - in this case, 4 hours.
		\end{block}
		\begin{block}{Option Analysis}
			\begin{itemize}
				\item \textbf{A - Incorrect:} MTD is total downtime
				\item \textbf{B - Correct:} RPO measures data loss tolerance
				\item \textbf{C - Incorrect:} RTO measures downtime tolerance
				\item \textbf{D - Incorrect:} BIA is the assessment process
			\end{itemize}
		\end{block}
	\end{frame}
	
	\begin{frame}{Q71: RTO and MTD}
		\begin{block}{Question}
			An application can tolerate four hours of downtime but must be recovered within six hours. Which metrics are correct?
		\end{block}
		\begin{enumerate}[(A)]
			\item RPO - 4h, RTO - 6h
			\item MTD - 4h, RTO - 6h
			\item MTD - 6h, RPO - 4h
			\item RTO - 4h, MTD - 6h
		\end{enumerate}
	\end{frame}
	
	\begin{frame}{Q71: Explanation}
		\begin{block}{Correct Answer: D}
			RTO is maximum acceptable downtime (4h). MTD is maximum total downtime before material impact (6h).
		\end{block}
		\begin{block}{Option Analysis}
			\begin{itemize}
				\item \textbf{A - Incorrect:} RPO is data loss, not downtime
				\item \textbf{B - Incorrect:} MTD should be 6h, not 4h
				\item \textbf{C - Incorrect:} MTD is 6h not 4h
				\item \textbf{D - Correct:} RTO=4h (downtime tolerance), MTD=6h (total before impact)
			\end{itemize}
		\end{block}
	\end{frame}
	
	\begin{frame}{Q72: Data Classification}
		\begin{block}{Question}
			Healthcare records and driver's licenses - most appropriate classification?
		\end{block}
		\begin{enumerate}[(A)]
			\item Sensitive (PII and PHI)
			\item Public
			\item PII
			\item PHI
		\end{enumerate}
	\end{frame}
	
	\begin{frame}{Q72: Explanation}
		\begin{block}{Correct Answer: A}
			Contains both PII and PHI, so classification should be Sensitive.
		\end{block}
		\begin{block}{Option Analysis}
			\begin{itemize}
				\item \textbf{A - Correct:} Sensitive covers both PII and PHI
				\item \textbf{B - Incorrect:} Public is for non-sensitive data
				\item \textbf{C - Incorrect:} PII alone doesn't cover health records
				\item \textbf{D - Incorrect:} PHI alone doesn't cover driver's licenses
			\end{itemize}
		\end{block}
	\end{frame}
	
	\begin{frame}{Q73: Data Destruction}
		\begin{block}{Question}
			Which stage of data life cycle includes disposal of data?
		\end{block}
		\begin{enumerate}[(A)]
			\item Destruction
			\item Archival
			\item Use
			\item Storage
		\end{enumerate}
	\end{frame}
	
	\begin{frame}{Q73: Explanation}
		\begin{block}{Correct Answer: A}
			Destruction stage includes disposal of data.
		\end{block}
		\begin{block}{Option Analysis}
			\begin{itemize}
				\item \textbf{A - Correct:} Destruction is disposal
				\item \textbf{B - Incorrect:} Archival is long-term storage
				\item \textbf{C - Incorrect:} Use is processing
				\item \textbf{D - Incorrect:} Storage is saving data
			\end{itemize}
		\end{block}
	\end{frame}
	
	\begin{frame}{Q74: Data Labeling Purpose}
		\begin{block}{Question}
			Data labeling is important for which of the following?
		\end{block}
		\begin{enumerate}[(A)]
			\item Data archival
			\item Data destruction
			\item DLP
			\item All of the above
		\end{enumerate}
	\end{frame}
	
	\begin{frame}{Q74: Explanation}
		\begin{block}{Correct Answer: D}
			Data labeling is important for archival, destruction, and DLP.
		\end{block}
		\begin{block}{Option Analysis}
			\begin{itemize}
				\item \textbf{A - Incorrect:} Archival alone not complete
				\item \textbf{B - Incorrect:} Destruction alone not complete
				\item \textbf{C - Incorrect:} DLP alone not complete
				\item \textbf{D - Correct:} All functions benefit from labeling
			\end{itemize}
		\end{block}
	\end{frame}
	
	\section{Domain 4: SDLC}
	
	\begin{frame}{Q75: SDLC Budget Phase}
		\begin{block}{Question}
			A project is budgeted in which phase of the SDLC?
		\end{block}
		\begin{enumerate}[(A)]
			\item Development
			\item Disposal
			\item Initiation
			\item Maintenance
		\end{enumerate}
	\end{frame}
	
	\begin{frame}{Q75: Explanation}
		\begin{block}{Correct Answer: C}
			Budgeting occurs in the initiation phase of SDLC.
		\end{block}
		\begin{block}{Option Analysis}
			\begin{itemize}
				\item \textbf{A - Incorrect:} Development is coding
				\item \textbf{B - Incorrect:} Disposal is retirement
				\item \textbf{C - Correct:} Initiation includes budgeting
				\item \textbf{D - Incorrect:} Maintenance is ongoing support
			\end{itemize}
		\end{block}
	\end{frame}
	
	\begin{frame}{Q76: User Training Phase}
		\begin{block}{Question}
			End user training is an integral part of which phase of the SDLC?
		\end{block}
		\begin{enumerate}[(A)]
			\item Development
			\item Disposal
			\item Initiation
			\item Implementation
		\end{enumerate}
	\end{frame}
	
	\begin{frame}{Q76: Explanation}
		\begin{block}{Correct Answer: D}
			End users are trained in the implementation phase.
		\end{block}
		\begin{block}{Option Analysis}
			\begin{itemize}
				\item \textbf{A - Incorrect:} Development is coding
				\item \textbf{B - Incorrect:} Disposal is retirement
				\item \textbf{C - Incorrect:} Initiation is planning
				\item \textbf{D - Correct:} Implementation includes training
			\end{itemize}
		\end{block}
	\end{frame}
	
	\begin{frame}{Q77: SDLC Final Stage}
		\begin{block}{Question}
			Which is the final stage of the SDLC?
		\end{block}
		\begin{enumerate}[(A)]
			\item Development
			\item Disposal
			\item Initiation
			\item Maintenance
		\end{enumerate}
	\end{frame}
	
	\begin{frame}{Q77: Explanation}
		\begin{block}{Correct Answer: B}
			Disposal is the final phase of the SDLC.
		\end{block}
		\begin{block}{Option Analysis}
			\begin{itemize}
				\item \textbf{A - Incorrect:} Development is early phase
				\item \textbf{B - Correct:} Disposal is final phase
				\item \textbf{C - Incorrect:} Initiation is first phase
				\item \textbf{D - Incorrect:} Maintenance is ongoing
			\end{itemize}
		\end{block}
	\end{frame}
	
	\begin{frame}{Q78: Project Risk vs SDLC Risk}
		\begin{block}{Question}
			Key difference between project risk and SDLC risk is:
		\end{block}
		\begin{enumerate}[(A)]
			\item Project risk relates to objectives, SDLC relates to time
			\item Project risk relates to objectives, SDLC relates to development
			\item Project risk relates to development, SDLC relates to objectives
			\item There is no difference
		\end{enumerate}
	\end{frame}
	
	\begin{frame}{Q78: Explanation}
		\begin{block}{Correct Answer: B}
			Project risk relates to project objectives; SDLC risk relates to development objectives.
		\end{block}
		\begin{block}{Option Analysis}
			\begin{itemize}
				\item \textbf{A - Incorrect:} SDLC risk is not just time
				\item \textbf{B - Correct:} Different focuses
				\item \textbf{C - Incorrect:} Reversed relationship
				\item \textbf{D - Incorrect:} They are different
			\end{itemize}
		\end{block}
	\end{frame}
	
	\begin{frame}{Q79: Certification Requirement}
		\begin{block}{Question}
			Which is a must to achieve certification?
		\end{block}
		\begin{enumerate}[(A)]
			\item A risk manager
			\item An internal audit team
			\item Policies
			\item An external audit firm
		\end{enumerate}
	\end{frame}
	
	\begin{frame}{Q79: Explanation}
		\begin{block}{Correct Answer: D}
			An external audit firm is required to certify a product.
		\end{block}
		\begin{block}{Option Analysis}
			\begin{itemize}
				\item \textbf{A - Incorrect:} Risk manager may not be required
				\item \textbf{B - Incorrect:} Internal audit cannot certify externally
				\item \textbf{C - Incorrect:} Policies alone don't certify
				\item \textbf{D - Correct:} External audit firm is required
			\end{itemize}
		\end{block}
	\end{frame}
	
	\begin{frame}{Q80: BYOD Policy}
		\begin{block}{Question}
			Practice of allowing employees to use personal laptops to access organizational resources is called:
		\end{block}
		\begin{enumerate}[(A)]
			\item Remote working
			\item BYOD
			\item Virtual computing
			\item Cloud computing
		\end{enumerate}
	\end{frame}
	
	\begin{frame}{Q80: Explanation}
		\begin{block}{Correct Answer: B}
			BYOD (Bring Your Own Device) allows personal devices to access organization resources.
		\end{block}
		\begin{block}{Option Analysis}
			\begin{itemize}
				\item \textbf{A - Incorrect:} Remote working is broader
				\item \textbf{B - Correct:} BYOD specifically allows personal devices
				\item \textbf{C - Incorrect:} Virtual computing is different
				\item \textbf{D - Incorrect:} Cloud computing is different
			\end{itemize}
		\end{block}
	\end{frame}
	
	\begin{frame}{Q81: Quantum Computing}
		\begin{block}{Question}
			Which technology has potential to break current encryption algorithms?
		\end{block}
		\begin{enumerate}[(A)]
			\item AI
			\item Blockchain
			\item Quantum computing
			\item IoT
		\end{enumerate}
	\end{frame}
	
	\begin{frame}{Q81: Explanation}
		\begin{block}{Correct Answer: C}
			Quantum computing has the potential to break current encryption algorithms.
		\end{block}
		\begin{block}{Option Analysis}
			\begin{itemize}
				\item \textbf{A - Incorrect:} AI doesn't break encryption
				\item \textbf{B - Incorrect:} Blockchain is distributed ledger
				\item \textbf{C - Correct:} Quantum computing breaks encryption
				\item \textbf{D - Incorrect:} IoT doesn't break encryption
			\end{itemize}
		\end{block}
	\end{frame}
	
	\section{Domain 4: Information Security}
	
	\begin{frame}{Q82: Encryption Keys}
		\begin{block}{Question}
			Which keys are required for asymmetric key encryption?
		\end{block}
		\begin{enumerate}[(A)]
			\item Only public key
			\item Only private key
			\item Any public or private key
			\item Public and private key pair
		\end{enumerate}
	\end{frame}
	
	\begin{frame}{Q82: Explanation}
		\begin{block}{Correct Answer: D}
			Asymmetric encryption requires both a public key and a private key.
		\end{block}
		\begin{block}{Option Analysis}
			\begin{itemize}
				\item \textbf{A - Incorrect:} Public key alone insufficient
				\item \textbf{B - Incorrect:} Private key alone insufficient
				\item \textbf{C - Incorrect:} Need both keys
				\item \textbf{D - Correct:} Both keys are required
			\end{itemize}
		\end{block}
	\end{frame}
	
	\begin{frame}{Q83: Symmetric Encryption}
		\begin{block}{Question}
			Which key is required for symmetric key encryption?
		\end{block}
		\begin{enumerate}[(A)]
			\item Public and private key
			\item A single key for encryption and decryption
			\item Only public key and hash
			\item Only private key and hash
		\end{enumerate}
	\end{frame}
	
	\begin{frame}{Q83: Explanation}
		\begin{block}{Correct Answer: B}
			Symmetric encryption uses a single key for both encryption and decryption.
		\end{block}
		\begin{block}{Option Analysis}
			\begin{itemize}
				\item \textbf{A - Incorrect:} That's asymmetric
				\item \textbf{B - Correct:} Single key for both
				\item \textbf{C - Incorrect:} Public key is asymmetric
				\item \textbf{D - Incorrect:} Private key is asymmetric
			\end{itemize}
		\end{block}
	\end{frame}
	
	\begin{frame}{Q84: Integrity Definition}
		\begin{block}{Question}
			Integrity of information means that:
		\end{block}
		\begin{enumerate}[(A)]
			\item Information is available to authorized users only
			\item Information is available at all times
			\item Information has not been tampered with
			\item All of these
		\end{enumerate}
	\end{frame}
	
	\begin{frame}{Q84: Explanation}
		\begin{block}{Correct Answer: C}
			Integrity means information has not been tampered with.
		\end{block}
		\begin{block}{Option Analysis}
			\begin{itemize}
				\item \textbf{A - Incorrect:} That's confidentiality
				\item \textbf{B - Incorrect:} That's availability
				\item \textbf{C - Correct:} Integrity is no tampering
				\item \textbf{D - Incorrect:} Only C is correct
			\end{itemize}
		\end{block}
	\end{frame}
	
	\begin{frame}{Q85: Confidentiality Definition}
		\begin{block}{Question}
			Confidentiality of information means that:
		\end{block}
		\begin{enumerate}[(A)]
			\item Information is available to authorized users only
			\item Information is available at all times
			\item Information has not been tampered with
			\item All of these
		\end{enumerate}
	\end{frame}
	
	\begin{frame}{Q85: Explanation}
		\begin{block}{Correct Answer: A}
			Confidentiality means information is available to authorized users only.
		\end{block}
		\begin{block}{Option Analysis}
			\begin{itemize}
				\item \textbf{A - Correct:} Confidentiality limits access
				\item \textbf{B - Incorrect:} That's availability
				\item \textbf{C - Incorrect:} That's integrity
				\item \textbf{D - Incorrect:} Only A is correct
			\end{itemize}
		\end{block}
	\end{frame}
	
	\begin{frame}{Q86: Availability Definition}
		\begin{block}{Question}
			Availability of information means that:
		\end{block}
		\begin{enumerate}[(A)]
			\item Information is available to authorized users only
			\item Information is available at all times to authorized users
			\item Information has not been tampered with
			\item All of these
		\end{enumerate}
	\end{frame}
	
	\begin{frame}{Q86: Explanation}
		\begin{block}{Correct Answer: B}
			Availability means information is accessible when needed by authorized users.
		\end{block}
		\begin{block}{Option Analysis}
			\begin{itemize}
				\item \textbf{A - Incorrect:} That's confidentiality
				\item \textbf{B - Correct:} Availability is timely access
				\item \textbf{C - Incorrect:} That's integrity
				\item \textbf{D - Incorrect:} Only B is correct
			\end{itemize}
		\end{block}
	\end{frame}
	
	\begin{frame}{Q87: Accountability}
		\begin{block}{Question}
			Which access control principle speaks to logging and monitoring of user activities?
		\end{block}
		\begin{enumerate}[(A)]
			\item Identification
			\item Authentication
			\item Authorization
			\item Accountability
		\end{enumerate}
	\end{frame}
	
	\begin{frame}{Q87: Explanation}
		\begin{block}{Correct Answer: D}
			Accountability deals with logging and monitoring of user activities.
		\end{block}
		\begin{block}{Option Analysis}
			\begin{itemize}
				\item \textbf{A - Incorrect:} Identification is unique ID
				\item \textbf{B - Incorrect:} Authentication is verification
				\item \textbf{C - Incorrect:} Authorization is permissions
				\item \textbf{D - Correct:} Accountability is logging/monitoring
			\end{itemize}
		\end{block}
	\end{frame}
	
	\begin{frame}{Q88: Identification}
		\begin{block}{Question}
			Which access control principle requires each user to have a separate unique name?
		\end{block}
		\begin{enumerate}[(A)]
			\item Identification
			\item Authentication
			\item Authorization
			\item Accountability
		\end{enumerate}
	\end{frame}
	
	\begin{frame}{Q88: Explanation}
		\begin{block}{Correct Answer: A}
			Identification requires all users to have a unique ID.
		\end{block}
		\begin{block}{Option Analysis}
			\begin{itemize}
				\item \textbf{A - Correct:} Identification requires unique ID
				\item \textbf{B - Incorrect:} Authentication verifies identity
				\item \textbf{C - Incorrect:} Authorization sets permissions
				\item \textbf{D - Incorrect:} Accountability logs activity
			\end{itemize}
		\end{block}
	\end{frame}
	
	\begin{frame}{Q89: Segregation of Duties}
		\begin{block}{Question}
			SOD principle ensures that:
		\end{block}
		\begin{enumerate}[(A)]
			\item All admin access to a single user
			\item Access to approve transactions is distributed
			\item Confidentiality is maintained
			\item Identification is in place
		\end{enumerate}
	\end{frame}
	
	\begin{frame}{Q89: Explanation}
		\begin{block}{Correct Answer: B}
			SOD ensures no single user has sole access to approve transactions.
		\end{block}
		\begin{block}{Option Analysis}
			\begin{itemize}
				\item \textbf{A - Incorrect:} SOD prevents single control
				\item \textbf{B - Correct:} SOD distributes authority
				\item \textbf{C - Incorrect:} Confidentiality is separate
				\item \textbf{D - Incorrect:} Identification is separate
			\end{itemize}
		\end{block}
	\end{frame}
	
	\begin{frame}{Q90: Hashing Characteristics}
		\begin{block}{Question}
			Which is true about hashing?
		\end{block}
		\begin{enumerate}[(A)]
			\item Provides unique output for each input
			\item Is non-reversible
			\item Helps maintain integrity
			\item All of these
		\end{enumerate}
	\end{frame}
	
	\begin{frame}{Q90: Explanation}
		\begin{block}{Correct Answer: D}
			All options are true for hashing - unique output, non-reversible, and integrity.
		\end{block}
		\begin{block}{Option Analysis}
			\begin{itemize}
				\item \textbf{A - True:} Unique output
				\item \textbf{B - True:} Non-reversible
				\item \textbf{C - True:} Integrity verification
				\item \textbf{D - Correct:} All are true
			\end{itemize}
		\end{block}
	\end{frame}
	
	\begin{frame}{Q91: Digital Signatures}
		\begin{block}{Question}
			Which is NOT provided by digital signatures?
		\end{block}
		\begin{enumerate}[(A)]
			\item Integrity
			\item Non-repudiation
			\item Authenticity
			\item Confidentiality
		\end{enumerate}
	\end{frame}
	
	\begin{frame}{Q91: Explanation}
		\begin{block}{Correct Answer: D}
			Digital signatures do NOT provide confidentiality on their own.
		\end{block}
		\begin{block}{Option Analysis}
			\begin{itemize}
				\item \textbf{A - Incorrect:} Integrity is provided
				\item \textbf{B - Incorrect:} Non-repudiation is provided
				\item \textbf{C - Incorrect:} Authenticity is provided
				\item \textbf{D - Correct:} Confidentiality is not provided
			\end{itemize}
		\end{block}
	\end{frame}
	
	\begin{frame}{Q92: Encryption and Confidentiality}
		\begin{block}{Question}
			Sender wants only receiver to decrypt data. Which should sender use?
		\end{block}
		\begin{enumerate}[(A)]
			\item Receiver's public key
			\item Sender's public key
			\item Receiver's private key
			\item Sender's private key
		\end{enumerate}
	\end{frame}
	
	\begin{frame}{Q92: Explanation}
		\begin{block}{Correct Answer: A}
			Use receiver's public key to encrypt so only receiver can decrypt with their private key.
		\end{block}
		\begin{block}{Option Analysis}
			\begin{itemize}
				\item \textbf{A - Correct:} Receiver's public key ensures only receiver decrypts
				\item \textbf{B - Incorrect:} Anyone with sender's public key
				\item \textbf{C - Incorrect:} Receiver's private key is not shared
				\item \textbf{D - Incorrect:} Sender's private key would allow anyone
			\end{itemize}
		\end{block}
	\end{frame}
	
	\begin{frame}{Q93: Signed and Encrypted}
		\begin{block}{Question}
			Which cryptographic mechanism provides confidentiality, integrity, and non-repudiation?
		\end{block}
		\begin{enumerate}[(A)]
			\item Signed, not encrypted
			\item Encrypted, not signed
			\item Signed, encrypted
			\item Signed, hashed
		\end{enumerate}
	\end{frame}
	
	\begin{frame}{Q93: Explanation}
		\begin{block}{Correct Answer: C}
			Signed and encrypted provides confidentiality, integrity, and non-repudiation.
		\end{block}
		\begin{block}{Option Analysis}
			\begin{itemize}
				\item \textbf{A - Incorrect:} No confidentiality
				\item \textbf{B - Incorrect:} No integrity/non-repudiation
				\item \textbf{C - Correct:} All three provided
				\item \textbf{D - Incorrect:} No confidentiality
			\end{itemize}
		\end{block}
	\end{frame}
	
	\begin{frame}{Q94: PKI SPOF}
		\begin{block}{Question}
			Which could be a SPOF in the PKI ecosystem?
		\end{block}
		\begin{enumerate}[(A)]
			\item Certificate
			\item CA
			\item Requester
			\item Browser
		\end{enumerate}
	\end{frame}
	
	\begin{frame}{Q94: Explanation}
		\begin{block}{Correct Answer: B}
			CA could be a SPOF if its private key is compromised.
		\end{block}
		\begin{block}{Option Analysis}
			\begin{itemize}
				\item \textbf{A - Incorrect:} Individual certificate isn't SPOF
				\item \textbf{B - Correct:} CA compromise affects all
				\item \textbf{C - Incorrect:} Requester isn't SPOF
				\item \textbf{D - Incorrect:} Browser isn't SPOF
			\end{itemize}
		\end{block}
	\end{frame}
	
	\begin{frame}{Q95: Security Awareness Primary Goal}
		\begin{block}{Question}
			Most important aspect of security awareness training is:
		\end{block}
		\begin{enumerate}[(A)]
			\item Train on latest threats
			\item Show compliance with audits
			\item Create security-conscious culture
			\item Upskill employees
		\end{enumerate}
	\end{frame}
	
	\begin{frame}{Q95: Explanation}
		\begin{block}{Correct Answer: C}
			Most important goal is to create a security-conscious culture.
		\end{block}
		\begin{block}{Option Analysis}
			\begin{itemize}
				\item \textbf{A - Incorrect:} Latest threats change
				\item \textbf{B - Incorrect:} Compliance is secondary
				\item \textbf{C - Correct:} Culture is primary
				\item \textbf{D - Incorrect:} Upskilling is secondary
			\end{itemize}
		\end{block}
	\end{frame}
	
	\begin{frame}{Q96: Data Privacy}
		\begin{block}{Question}
			Data privacy should be:
		\end{block}
		\begin{enumerate}[(A)]
			\item Each individual's right
			\item Supported by security controls
			\item Embedded into system design
			\item All of these
		\end{enumerate}
	\end{frame}
	
	\begin{frame}{Q96: Explanation}
		\begin{block}{Correct Answer: D}
			Data privacy should be each individual's right, supported by controls, and embedded in design.
		\end{block}
		\begin{block}{Option Analysis}
			\begin{itemize}
				\item \textbf{A - Incorrect:} Right alone not enough
				\item \textbf{B - Incorrect:} Controls alone not enough
				\item \textbf{C - Incorrect:} Design alone not enough
				\item \textbf{D - Correct:} All aspects are needed
			\end{itemize}
		\end{block}
	\end{frame}
	
	\begin{frame}{Q97: Data Security vs Privacy}
		\begin{block}{Question}
			Which statement is true about data security and data privacy?
		\end{block}
		\begin{enumerate}[(A)]
			\item They are the same thing
			\item Security is protecting data; privacy is protecting personal information
			\item Privacy is protecting data; security is protecting personal information
			\item Not important for organizations
		\end{enumerate}
	\end{frame}
	
	\begin{frame}{Q97: Explanation}
		\begin{block}{Correct Answer: B}
			Data security protects data from unauthorized access; data privacy protects personal information from misuse.
		\end{block}
		\begin{block}{Option Analysis}
			\begin{itemize}
				\item \textbf{A - Incorrect:} They are different
				\item \textbf{B - Correct:} Security vs privacy distinction
				\item \textbf{C - Incorrect:} Reversed
				\item \textbf{D - Incorrect:} Both are important
			\end{itemize}
		\end{block}
	\end{frame}
	
	\begin{frame}{Q98: Privacy Compliance}
		\begin{block}{Question}
			Company collecting customer data for marketing. Which measure ensures compliance with privacy regulations?
		\end{block}
		\begin{enumerate}[(A)]
			\item Use encryption
			\item Obtain consent
			\item Implement firewalls
			\item Train employees
		\end{enumerate}
	\end{frame}
	
	\begin{frame}{Q98: Explanation}
		\begin{block}{Correct Answer: B}
			Obtaining consent is critical for compliance with GDPR, CCPA, and other privacy regulations.
		\end{block}
		\begin{block}{Option Analysis}
			\begin{itemize}
				\item \textbf{A - Incorrect:} Encryption is security, not privacy
				\item \textbf{B - Correct:} Consent is required
				\item \textbf{C - Incorrect:} Firewalls are security
				\item \textbf{D - Incorrect:} Training is important but not primary
			\end{itemize}
		\end{block}
	\end{frame}
	
	\begin{frame}{Q99: Vulnerability Management KPI}
		\begin{block}{Question}
			Which KPI is most appropriate to measure vulnerability management program effectiveness?
		\end{block}
		\begin{enumerate}[(A)]
			\item Number of vulnerabilities found and remediated
			\item Number of security incidents reported
			\item Percentage of employees trained
			\item Average incident response time
		\end{enumerate}
	\end{frame}
	
	\begin{frame}{Q99: Explanation}
		\begin{block}{Correct Answer: A}
			Number of vulnerabilities found and remediated measures vulnerability management effectiveness.
		\end{block}
		\begin{block}{Option Analysis}
			\begin{itemize}
				\item \textbf{A - Correct:} Direct measure of VM program
				\item \textbf{B - Incorrect:} Measures incident response
				\item \textbf{C - Incorrect:} Measures training program
				\item \textbf{D - Incorrect:} Measures IR effectiveness
			\end{itemize}
		\end{block}
	\end{frame}
	
	\begin{frame}{Q100: Security Program KPI}
		\begin{block}{Question}
			Which KPI is most useful to evaluate security program effectiveness over time?
		\end{block}
		\begin{enumerate}[(A)]
			\item Number of employees trained
			\item Number of security incidents reported
			\item Average detection and response time
			\item Percentage reduction in data breaches
		\end{enumerate}
	\end{frame}
	
	\begin{frame}{Q100: Explanation}
		\begin{block}{Correct Answer: D}
			Percentage reduction in data breaches directly measures program impact over time.
		\end{block}
		\begin{block}{Option Analysis}
			\begin{itemize}
				\item \textbf{A - Incorrect:} Training completion not a direct measure
				\item \textbf{B - Incorrect:} Incidents may change
				\item \textbf{C - Incorrect:} Response time is operational
				\item \textbf{D - Correct:} Direct measure of program success
			\end{itemize}
		\end{block}
	\end{frame}
	
	% ================ FINAL SLIDES ================
	
	\begin{frame}{CRISC Key Takeaways}
		\begin{block}{Essential Concepts}
			\begin{itemize}
				\item Governance provides oversight and direction from Board
				\item Risk identification is the first step - "You can't protect what you don't know"
				\item Risk Appetite = willingness to accept risk
				\item Risk Tolerance = acceptable variation from appetite
				\item Risk Capacity = maximum risk before existence threatened
				\item RPO is backward-looking; RTO is forward-looking
			\end{itemize}
		\end{block}
	\end{frame}
	
	\begin{frame}{CRISC Key Takeaways (continued)}
		\begin{block}{More Essential Concepts}
			\begin{itemize}
				\item 3LoD: First = Operations, Second = Oversight, Third = Audit
				\item Residual Risk = Inherent Risk - Controls
				\item Risk Response: Mitigate, Accept, Transfer, Avoid
				\item SMART: Specific, Measurable, Attainable, Relevant, Timely
				\item CIA: Confidentiality, Integrity, Availability
				\item Digital signatures: Integrity, Non-repudiation, Authenticity (NOT confidentiality)
			\end{itemize}
		\end{block}
	\end{frame}
	
	\begin{frame}{Final Advice}
		\begin{block}{Exam Preparation}
			\centering
			\vspace{0.3cm}
			\textbf{Remember:}
			\begin{itemize}
				\item Understand the \textit{why} behind each concept
				\item Know control categories and their differences
				\item Understand recovery objectives (RPO, RTO, MTD)
				\item Know risk appetite vs tolerance vs capacity
				\item Practice scenario-based questions
			\end{itemize}
			\vspace{0.5cm}
			\textbf{Good Luck on Your CRISC Exam!}
		\end{block}
	\end{frame}
	
\end{document}